%% ============================================================
%% 片段：sB-notation   对应原文 PDF 第 153--160 页（Appendix B）
%% 译者：lead   审校状态：初稿
%% ============================================================

% ==== 原文 Appendix B：PDF 第 153 页（印刷第 152 页）====
\section{记号索引}
\label{sec:B}

本索引按原文的四组编排：I.\ 规范场论；II.\ 超 K\"ahler 结构；III.\ 群论；IV.\ 格。
左列为符号，右列为含义（原文 Appendix B 共 91 条，本译本逐条译出，符号一仍其旧）。

\begin{yizhu}
本索引中两个 $\mathrm{N}$ 形符号极易混淆：${\mathfrak N}$（原文 \verb|\eurm N|，Euler Roman）是\textbf{瞬子数}，
而 $\EUN$（花体 $\mathcal N$）是抛物子群 $\EUP$ 的\textbf{幂幺根}；同理 ${\mathfrak D}$（仿射外尔房）与
$\CF$、$\cD$ 无关。原文用两套字体（Euler Roman 与 Euler Script）区分，中译本以
$\mathfrak{N}$／$\mathfrak{D}$ 与 $\mathcal{N}$／$\mathcal{D}$ 对应，改稿时请勿互换。
\end{yizhu}

\begingroup\small
\begin{longtable}{@{}l@{\hskip 1.2em}>{\raggedright\arraybackslash}p{0.64\linewidth}@{}}
\multicolumn{2}{@{}l}{\textbf{I.\ 规范场论}}\\
\midrule
$M$ & 四维流形\\
$W$ & 三维流形\\
$C$ & 我们做几何朗兰兹纲领的黎曼面\\
$\Sigma$ & 我们紧化到的黎曼面\\
$G$ & 紧致规范群，通常取单群\\
$\bar G$ & $G$ 的单连通覆盖\\
$E$ & 一个 $G$-丛\\
$A$ & 规范场；$E$ 上的联络\\
$d_A$ & 规范协变外微分 $d+[A,\,\cdot\,]$\\
$\phi$ & ${\rm ad}(E)$-值 1-形式；希格斯场\\
$\varphi$ & $\phi$ 的 $(1,0)$ 部分\\
${\mathcal A}$ & 复值联络 $A+i\phi$\\
$V$ & ${\mathcal A}$ 的单值性\\
$U$ & $V$ 的轨道闭包中的半单元素\\
$\tau$ & 规范耦合参数 $\theta/2\pi+4\pi i/e^2$\\
$\mathfrak{N}$ & 瞬子数\\
$\mathfrak{m}$ & 磁通量；$\mathbb{T}$-丛的特征类\\
$\xi$ & 阻碍 $G$-丛 $E$ 提升为 $\bar G$-丛的特征类\\
$D$ & 面算符的支撑\\
$L$ & 线算符的支撑\\
$(\alpha,\beta,\gamma,\eta)$ & 面算符的参数\\
${\bf e}_0$ & 离散电通量\\
${\bf m}_0$ & 离散磁通量\\
$\MH$ & 希格斯丛的模空间\\
$\M$ & $G$-丛的模空间\\
$\MH(\alpha,\beta,\gamma)$ & 带分歧希格斯丛的模空间\\
$\M(\alpha)$ & 抛物丛的模空间\\
$\EUBB$ & 希钦纤维化的底空间\\
$\CMF$ & 希钦纤维化的纤维\\
${\mathcal B}$ & 一个膜（brane）\\
${\mathcal L}$ & 膜的 Chan--Paton 线丛\\
$\EUG$ & 规范变换群\\
$\star$ & 霍奇星算符\\
$\Gamma$ & 对偶群\\
\midrule
\multicolumn{2}{@{}l}{\textbf{II.\ 超 K\"ahler 结构}}\\
\midrule
$X$ & 一般的超 K\"ahler 流形\\
$I,J,K$ & $X$ 上的复结构\\
$\omega_I,\omega_J,\omega_K$ & $X$ 上的 K\"ahler 结构\\
$\Omega_I,\Omega_J,\Omega_K$ & $X$ 上的全纯 2-形式\\
$\vec\mu=(\mu_I,\mu_J,\mu_K)$ & 超 K\"ahler 矩映射\\
\midrule
\multicolumn{2}{@{}l}{\textbf{III.\ 群论}}\\
\midrule
$G$ & 紧致李群，通常取单群\\
$\bar G$ & $G$ 的万有覆盖\\
$G_{\rm ad}$ & $G$ 的伴随形式\\
$^L\neg G$ & 对偶群\\
$\mathbb{T}$ & $G$ 的极大环面\\
$\mathbb{L}$ & $G$ 的列维子群\\
$G_\C,\,\mathbb{T}_\C,$ 等 & 复化\\
$\EUB$ & $G_\C$ 的博雷尔子群\\
$\EUP$ & $G_\C$ 的抛物子群\\
$\EUN$ & $\EUP$ 的幂幺根\\
$\mathfrak g$ & $G$ 的李代数\\
$\mathfrak t$ & $\mathbb{T}$ 的李代数\\
$\mathfrak t^\vee$ & $\mathfrak t$ 的对偶\\
$\mathfrak b$ & $\EUB$ 的李代数\\
$\mathfrak p$ & $\EUP$ 的李代数\\
$\mathfrak n$ & $\EUN$ 的李代数\\
${\mathcal W}$ & 外尔群\\
${\mathcal W}_{\rm aff}$ & 仿射外尔群\\
$\mathfrak{D}$ & 仿射外尔房\\
$B_{\rm aff}$ & 仿射辫群\\
${\mathcal Z}(G)$ & $G$ 的中心\\
$\pi_1(G)$ & $G$ 的基本群\\
$\mathfrak C$ & $G_\C$ 中的共轭类\\
$\mathfrak c$ & $\mathfrak g_\C$ 中的共轭类\\
$\mathfrak C_{\mathbb{L}}$ & 与列维子群 $\mathbb{L}$ 相联系的 Richardson 共轭类\\
\midrule
\multicolumn{2}{@{}l}{\textbf{IV.\ 格}}\\
\midrule
$\Lambda_\rt$ & 根格\\
$\Lambda_\cort$ & 余根格\\
$\Lambda_\wt$ & 权格\\
$\Lambda_\cowt$ & 余权格\\
$\Lambda_\char$ & 特征格\\
$\Lambda_\cochar$ & 余特征格\\
\end{longtable}
\endgroup
